% from ega4-21b-fr.tex
\oldpage[IV-4]{273}
donc $\operatorname{cyc}(D|U)=\operatorname{cyc}(D)|U$, et les homomorphismes $\operatorname{cyc}:\operatorname{Div}^+(U)\to\mathfrak{Z}^1(U)$ définissent donc un homomorphisme de faisceaux de monoïdes \sref{IV.21.6.4.2}, d'où l'homomorphisme cherché \sref{IV.21.6.4.1} de faisceaux de groupes.

On a $\Supp(\operatorname{cyc}(D))\subset\Supp(D)$ pour tout diviseur $D$ et
\begin{env}[21.6.6.2]
\label{IV.21.6.6.2}
\[
  \Supp(\operatorname{cyc}(D))=\Supp(D)\quad\text{lorsque}\quad D\geq0;
\]
\end{env}
on a déjà vu en effet la seconde relation \sref{IV.21.2.12};
lorsque $D$ est quelconque, la relation $D_x=0$ entraîne l'existence d'un voisinage ouvert $U$ de $x$ tel que $D|U=0$, d'où $\operatorname{cyc}(D)|U=0$, ce qui prouve notre assertion.

\begin{env}[21.6.7]
\EGAReaderTarget{ega4.c04.statement.21-6-7}{IV.21.6.7}
L'homomorphisme \sref{IV.21.6.4.1} donne, par passage aux groupes de sections au-dessus de $X$, un \emph{homomorphisme croissant de groupes ordonnés} $\operatorname{Div}(X)\to\mathfrak{Z}^1(X)$, que nous noterons encore $D\mapsto\operatorname{cyc}(D)$;
le nombre $\operatorname{mult}_x(\operatorname{cyc}(D))$ pour $x\in X^{(1)}$ se note encore $\operatorname{mult}_x(D)$ ou $\operatorname{mult}_{\overline{\{x\}}}(D)$ et s'appelle \emph{multiplicité du diviseur $D$ au point $x$}, ou \emph{multiplicité du cycle premier $\overline{\{x\}}$ dans $D$};
c'est un entier positif ou négatif, égal comme on l'a vu à $\operatorname{long}(\sh{O}_{Y(D),x})$ lorsque $D$ est un diviseur \emph{positif};
on a donc par définition
\[
\label{IV.21.6.7.1}
  \operatorname{cyc}(D)=\sum_{x\in X^{(1)}}\operatorname{mult}_x(D)\cdot\overline{\{x\}},
  \tag{21.6.7.1}
\]
et en vertu du fait que $D\mapsto\operatorname{cyc}(D)$ est un homomorphisme, on a
\[
\label{IV.21.6.7.2}
  \operatorname{mult}_x(-D)=-\operatorname{mult}_x(D),\qquad
  \operatorname{mult}_x(D+D')=\operatorname{mult}_x(D)+\operatorname{mult}_x(D')
  \tag{21.6.7.2}
\]
pour deux diviseurs quelconques $D$, $D'$.

Pour toute fonction méromorphe régulière $f$ sur $X$, on pose en particulier, pour $x\in X^{(1)}$
\[
\label{IV.21.6.7.3}
  o_x(f)=\operatorname{mult}_x(\operatorname{div}(f))
  \tag{21.6.7.3}
\]
et on dit que $o_x(f)$ (entier positif ou négatif) est \emph{l'ordre de $f$ au point} $x\in X^{(1)}$.
Si $f_x\in\sh{O}_{X,x}$ (élément régulier de cet anneau local par hypothèse), on a donc
\[
\label{IV.21.6.7.4}
  o_x(f)=\operatorname{long}(\sh{O}_{X,x}/f_x\sh{O}_{X,x});
  \tag{21.6.7.4}
\]
pour deux fonctions méromorphes régulières $f$, $f'$ sur $X$, on a
\[
\label{IV.21.6.7.5}
  o_x(ff')=o_x(f)+o_x(f'),\qquad o_x(f^{-1})=-o_x(f)
  \tag{21.6.7.5}
\]
pour tout $x\in X^{(1)}$.
Les cycles $1$-codimensionnels
\[
  Z^+(f)=\sum_{x\in X^{(1)}}(o_x(f))^+\cdot\overline{\{x\}},\qquad
  Z^-(f)=\sum_{x\in X^{(1)}}(o_x(f))^-\cdot\overline{\{x\}}
\]
sont appelés respectivement le \emph{cycle des zéros} et le \emph{cycle des pôles} (ou \emph{cycle polaire}) de $f$, et l'on a évidemment $\operatorname{cyc}(\operatorname{div}(f))=Z^+(f)-Z^-(f)$.
Les cycles $1$-codimensionnels de la forme $\operatorname{cyc}(\operatorname{div}(f))$ sont appelés \emph{principaux} (ou \emph{linéairement équivalents à zéro}) et forment un sous-groupe $\mathfrak{Z}^1_\mathrm{princ}(X)$ de $\mathfrak{Z}^1(X)$.
Les sections au-dessus de $X$ de l'image $c(\sh{D}iv_X)\subset\mathscr{Z}_X^1$ sont
